\subsection{Setup}\label{sec:setup}

\textbf{Traffic.} We use all nine devices of N-BaIoT~\cite{nbaiot}: two doorbells (Danmini, Ennio), a
thermostat (Ecobee), a baby monitor (Philips), four security cameras (Provision PT-737E and PT-838,
SimpleHome XCS7-1002 and -1003) and a webcam (Samsung), each with its real benign capture and its real
Mirai and BASHLITE attack traffic (Ennio and Samsung have BASHLITE only). Each row is the 115 Kitsune
statistics of one packet~\cite{kitsune}. The first 20\% of each benign capture calibrates the device's
detector; the rest streams in capture order, alternate rows for learning and for measuring false alarms.
Each attack type is split into a pool the poisoner draws from and a disjoint held-out set on which
detection is measured. The shared 16-dimensional whitened basis keeps \ExplainedVar\ of the within-device variance.

\textbf{Fleet.} The main fleet has the nine devices, one each; peers are all devices of other vendors.
For the monoculture experiment a twelve-device fleet has three instances each of four models; instances
differ in calibration subsample, position in the capture and random seed.

\textbf{Honest drift.} We apply drift to the raw Kitsune statistics with physical meaning: packet rate
$\times g$ (counts scale by $g$, inter-packet gaps by $1/g$), $b$ bytes of per-packet overhead in both
directions (a tunnel or a new record layer), and network jitter $\times h$. Attack traffic gets the
overhead and jitter (same network path) but not the benign rate change. Kitsune's log-scaled
statistics barely move under small changes (packet rate $\times 8$ raises a static detector's false
alarms only to 22\%), so drift strength was fixed by its effect on the static detector alone, never by
how any adaptive rule behaves. \emph{Strong} common drift ($g=64$, $b=240$, $h=11.4$, ramped over cycles 40
to 160) takes the static detector to \StrongStaticFPR\ false alarms; \emph{moderate} ($g=16$, $b=160$,
$h=5$) to about a third. Device-only drift ($g=32$, $b=200$, $h=7.6$, cycles 60 to 120) models one
device's own large change. These are stress tests of a large shared change; Section~\ref{sec:cesnet}
uses real drift.

\textbf{Attacker.} The poisoner targets the victim's scan traffic (Mirai scan, or BASHLITE scan on the
two devices without Mirai), the attack class closest to benign traffic, from cycle 40 on, with share
$p\in\{0.1,0.2,0.3,0.4\}$ of the victim's learning batch. Modes: \emph{boundary} (real benign points
pulled toward real attack points until just inside the alarm radius), \emph{core} (just inside the core
radius), \emph{mimicry} (re-centred copies of benign traffic shifted toward the attack, which looks like
drift to the device), \emph{multi-direction}, \emph{consensus-aligned} and \emph{pulsing}.

\textbf{Baselines.} All rules share the same detector and the same local proposal, so differences come
from how the proposal is applied. \emph{Static} never adapts. \emph{Naive} applies a mean-and-radius
proposal over all accepted traffic. \emph{Single-device} applies our signature-free proposal in full.
\emph{Clipping} caps each update's norm (robust-aggregation-style clipping with no memory).
\emph{Frozen reserve} learns only from points a frozen copy of the calibrated detector accepts
(WARDEN-style screening~\cite{warden}). \emph{Alarm-rate switch-off} skips updates while the batch alarm
rate exceeds five times the clean rate (Kloft and Laskov~\cite{kloft}). \emph{Outlier-ratio gate} skips
updates while the alarm rate exceeds three times its own moving average (DriftXpert style~\cite{driftxpert}).
\emph{TOA} is our earlier rule that rations movement toward an \emph{oracle} attack signature.
\emph{Same-model consensus} applies the element-wise median of same-model devices' proposals and discards
the rest (CIoTA style~\cite{ciota}). The gates are our re-implementations of the gating idea in each
paper, not the full published systems.

\textbf{Metrics.} A victim is \emph{compromised} when fewer than half of held-out target attack packets
are flagged; time to compromise (TTC) is the first such cycle. False-alarm rate (FPR) is measured on
held-out benign traffic each cycle, on the victim and on the rest of the fleet. All results are over 10
seeds (95\% confidence intervals where shown); a run is 240 cycles of 200 learning points per device.
